%! TEX root = standar_peraturan.tex
% the main TeX file which is intended to compile, :VimtexReload after adjustment
% :h vimtex-tex-root
\documentclass[../main.tex]{subfiles}
% \includeonlyframes{current}
% \setbeameroption{hide notes} % Only slides
% \setbeameroption{show only notes} % Only notes
% \setbeameroption{show notes on second screen=right} % Both

\begin{document}

% Slide 1
\begin{frame}[mytitle=center, c, mycolor=digiPH_cream]
	\titlepage
	\begin{itemize}
		\item \textbf{Fokus Pembelajaran:} Menghubungkan regulasi, kondisi eksisting, dan keputusan rancangan.
		\item \textbf{Ilustrasi:} Gambar tapak dengan batas persil, bangunan, jalan, dan zona pembatas pembangunan.
	\end{itemize}
\end{frame}

% Slide 2
\begin{frame}[mytitle=standard, t, mycolor=digiPH_gray]
	\frametitle{Konsep Dasar Standar dan Peraturan Tapak}
	\framesubtitle{Pengertian dan Fungsi Regulasi}
	\begin{itemize}
		\item Pengertian standar, regulasi tata ruang, ketentuan bangunan, dan analisis regulatif tapak.
		\item Perbedaan antara kondisi fisik yang ada dan batasan hukum yang mengatur pembangunan.
		\item Fungsi regulasi sebagai dasar menentukan kelayakan, kapasitas, dan arah pengembangan tapak.
		\item Hubungan regulasi dengan keselamatan, kenyamanan, aksesibilitas, dan keberlanjutan.
	\end{itemize}
\end{frame}

% Slide 3 - Part 1
\begin{frame}[mytitle=standard, t, mycolor=white]
	\frametitle{Landasan Hukum Penataan Tapak}
	\framesubtitle{Undang-Undang dan Peraturan Pemerintah}
	\begin{itemize}
		\item UU Nomor 26 Tahun 2007 tentang Penataan Ruang beserta perubahannya melalui UU Nomor 6 Tahun 2023.
		\item PP Nomor 21 Tahun 2021 tentang Penyelenggaraan Penataan Ruang.
		\item PP Nomor 16 Tahun 2021 tentang Peraturan Pelaksanaan UU Bangunan Gedung.
	\end{itemize}
\end{frame}

% Slide 3 - Part 2
\begin{frame}[mytitle=standard, t, mycolor=white]
	\frametitle{Landasan Hukum Penataan Tapak}
	\framesubtitle{Peraturan Menteri dan Modul Pengayaan}
	\begin{itemize}
		\item Permen PUPR Nomor 6 Tahun 2007 tentang Pedoman Umum Penyusunan RTBL. \textit{(Status tercatat berlaku pada basis data BPK)}
		\item Permen ATR/Kepala BPN Nomor 11 Tahun 2021 sebagaimana diubah dengan Permen Nomor 6 Tahun 2026.
		\item Modul 4 Rencana Umum dan Panduan Rancang, Kementerian PUPR (2016), sebagai bahan pengayaan.
		\item \textbf{Catatan:} Tekankan fungsi setiap regulasi dan perlunya memeriksa naskah resmi serta ketentuan daerah yang berlaku.
	\end{itemize}
\end{frame}

% Slide 4
\begin{frame}[mytitle=imageplus, c, mycolor=digiPH_gray]
	\frametitle{Hierarki dan Dokumen Tata Ruang}
	\begin{itemize}
		\item Peran RTRW dalam menentukan arah dan struktur pemanfaatan ruang wilayah.
		\item Peran RDTR dan peraturan zonasi dalam memberikan ketentuan yang lebih terperinci.
		\item Identifikasi ketentuan daerah, rencana tata bangunan, dan informasi persil yang relevan.
		\item Pemeriksaan lokasi, fungsi kegiatan, zona, serta ketentuan yang berlaku pada tapak studi.
		\item Prinsip verifikasi: gunakan dokumen resmi dan versi yang berlaku, bukan hanya peta ilustratif.
	\end{itemize}
\end{frame}

% Slide 5
\begin{frame}[mytitle=standard, t, mycolor=white]
	\frametitle{Analisis Peruntukan dan Kegiatan Lahan}
	\begin{itemize}
		\item Identifikasi zona dan fungsi lahan yang diizinkan, diizinkan bersyarat, atau tidak diizinkan sesuai peraturan setempat.
		\item Perbandingan fungsi yang direncanakan dengan fungsi yang diperbolehkan.
		\item Pengaruh peruntukan terhadap jenis bangunan, kegiatan, kebutuhan parkir, dan dampak lingkungan.
		\item Pemetaan zona yang sesuai, zona bersyarat, serta area yang berpotensi mengalami konflik pemanfaatan.
	\end{itemize}
\end{frame}

% Slide 6
\begin{frame}[mytitle=standard, c, mycolor=digiPH_cream]
	\frametitle{Batas Persil dan Garis Sempadan}
	\begin{itemize}
		\item Penentuan batas persil berdasarkan dokumen yang dapat dipertanggungjawabkan.
		\item Pengertian dan fungsi garis sempadan bangunan serta sempadan jalan atau sempadan lain yang relevan.
		\item Identifikasi batas pembangunan di sisi depan, samping, dan belakang sesuai ketentuan setempat.
		\item Pemetaan area yang dapat dibangun dan area yang harus bebas dari bangunan.
		\item Implikasi terhadap bentuk, posisi, dan orientasi massa bangunan.
	\end{itemize}
\end{frame}

% Slide 7
\begin{frame}[mytitle=standard, t, mycolor=white]
	\frametitle{Koefisien Dasar Bangunan (KDB)}
	\begin{itemize}
		\item Pengertian KDB sebagai perbandingan luas lantai dasar bangunan yang diperhitungkan terhadap luas persil.
		\item \textbf{Rumus:} KDB (\%) = (Luas dasar bangunan $\div$ Luas persil) $\times$ 100\%
		\item Contoh ilustratif: Persil 1.000 m² dengan KDB maksimum 60\% menghasilkan batas luas dasar 600 m².
		\item Penerapan pada alternatif tapak bangunan dan evaluasi ruang terbuka.
		\item \textit{Penegasan bahwa angka 60\% hanya contoh perhitungan, bukan standar nasional yang berlaku seragam.}
	\end{itemize}
\end{frame}

% Slide 8
\begin{frame}[mytitle=standard, t, mycolor=digiPH_gray]
	\frametitle{Koefisien Lantai Bangunan (KLB)}
	\begin{itemize}
		\item Pengertian KLB sebagai perbandingan luas total lantai yang diperhitungkan terhadap luas persil.
		\item \textbf{Rumus:} KLB = Luas total lantai bangunan $\div$ Luas persil
		\item Contoh ilustratif: Persil 1.000 m² dengan KLB 2,0 menghasilkan batas luas total lantai 2.000 m².
		\item Hubungan KLB dengan jumlah lantai, ketinggian, kepadatan, dan bentuk massa.
		\item Pemeriksaan terhadap definisi luas lantai serta batas ketinggian yang berlaku setempat.
	\end{itemize}
\end{frame}

% Slide 9
\begin{frame}[mytitle=standard, c, mycolor=white]
	\frametitle{KDH dan KTB dalam Pengembangan Tapak}
	\begin{itemize}
		\item \textbf{KDH:} Koefisien yang mengatur proporsi ruang hijau sesuai definisi dalam peraturan yang berlaku.
		\item \textbf{KTB:} Koefisien yang mengatur luas tapak basemen yang diperhitungkan sesuai ketentuan setempat.
		\item Hubungan ruang hijau, resapan air, area terbangun, dan ruang bawah tanah.
		\item Ilustrasi pembagian persil untuk bangunan, vegetasi, dan area permeabel.
		\item Verifikasi rumus, definisi, serta nilai minimum atau maksimum berdasarkan RDTR dan peraturan daerah.
	\end{itemize}
\end{frame}

% Slide 10
\begin{frame}[mytitle=standard, t, mycolor=digiPH_cream]
	\frametitle{Tata Bangunan dan Ketinggian}
	\begin{itemize}
		\item Ketentuan ketinggian, jumlah lantai, jarak antarmassa, orientasi, dan bentuk bangunan.
		\item Hubungan tata bangunan dengan garis sempadan dan intensitas pemanfaatan lahan.
		\item Pengaruh pencahayaan alami, penghawaan, privasi, bayangan, dan karakter visual kawasan.
		\item Pengecekan ketentuan khusus, termasuk kawasan bersejarah atau pembatasan keselamatan penerbangan apabila relevan.
		\item Penyusunan alternatif massa yang memenuhi ketentuan dan merespons lingkungan.
	\end{itemize}
\end{frame}

% Slide 11
\begin{frame}[mytitle=standard, t, mycolor=white]
	\frametitle{Ruang Terbuka Hijau dan Kualitas Lingkungan}
	\begin{itemize}
		\item Fungsi ruang terbuka hijau untuk resapan, pengurangan panas, kenyamanan, dan kualitas ekologis.
		\item Integrasi vegetasi eksisting, pohon pelindung, jalur hijau, serta ruang terbuka dalam rancangan.
		\item Pasal 29 UU Nomor 26 Tahun 2007 menetapkan proporsi RTH paling sedikit 30\% dari luas wilayah kota, termasuk RTH publik paling sedikit 20\%.
		\item \textbf{Catatan:} Angka tersebut berlaku pada skala wilayah kota dan bukan otomatis persyaratan minimum setiap persil.
		\item Ketentuan RTH pada tapak harus diperiksa kembali melalui regulasi lokal yang berlaku.
	\end{itemize}
\end{frame}

% Slide 12
\begin{frame}[mytitle=center, c, mycolor=digiPH_gray]
	\frametitle{Sirkulasi, Aksesibilitas, dan Keselamatan}
	\begin{itemize}
		\item Analisis akses kendaraan, pejalan kaki, pesepeda, kendaraan darurat, dan akses servis.
		\item Hubungan pintu masuk, parkir, jalur internal, serta koneksi dengan jaringan jalan.
		\item Prinsip aksesibilitas bagi penyandang disabilitas dan pengguna dengan kebutuhan beragam.
		\item Pemeriksaan konflik kendaraan dan pejalan kaki, visibilitas, keselamatan, dan kelancaran sirkulasi.
		\item Pemetaan jalur utama, titik konflik, dan alternatif akses yang memenuhi ketentuan.
	\end{itemize}
\end{frame}

% Slide 13
\begin{frame}[mytitle=standard, t, mycolor=white]
	\frametitle{Prasarana dan Utilitas Lingkungan}
	\begin{itemize}
		\item Sistem drainase, penyediaan air bersih, air limbah, persampahan, listrik, dan telekomunikasi.
		\item Identifikasi jaringan eksisting, jalur pelayanan, ruang pemeliharaan, dan potensi konflik dengan bangunan.
		\item Analisis limpasan air hujan, titik genangan, elevasi, serta hubungan dengan topografi.
		\item Pemetaan utilitas dan area perlindungan sesuai persyaratan teknis terkait.
		\item Integrasi utilitas dengan ruang hijau dan sirkulasi agar tidak mengurangi fungsi tapak.
	\end{itemize}
\end{frame}

% Slide 14
\begin{frame}[mytitle=standard, t, mycolor=digiPH_cream]
	\frametitle{Integrasi Regulasi dengan Kondisi Eksisting}
	\begin{itemize}
		\item Overlay batas persil, zonasi, garis sempadan, topografi, akses, vegetasi, drainase, dan bangunan eksisting.
		\item Pembedaan kendala regulatif dengan kendala fisik, ekologis, dan operasional.
		\item \textbf{Identifikasi peluang:} Akses yang baik, bentuk persil efisien, dan ruang hijau yang dapat dipertahankan.
		\item \textbf{Identifikasi risiko:} Zona terlarang, lereng bermasalah, genangan, atau konflik utilitas.
		\item Penyusunan peta gabungan sebagai dasar evaluasi alternatif desain.
	\end{itemize}
\end{frame}

% Slide 15
\begin{frame}[mytitle=standard, c, mycolor=white]
	\frametitle{Penyusunan Peta Kendala dan Potensi Regulatif}
	\begin{itemize}
		\item \textbf{Tahapan:} pengumpulan data, verifikasi dokumen, georeferensi, digitasi, analisis spasial, dan pemeriksaan lapangan.
		\item Kategori peta: area yang tidak dapat dibangun, area bersyarat, dan area potensial sesuai ketentuan.
		\item Penggunaan legenda, skala, arah utara, sumber data, tanggal, dan sistem koordinat yang jelas.
		\item Pencantuman nilai KDB, KLB, KDH, KTB, sempadan, dan ketinggian yang sudah diverifikasi.
		\item Hasil akhir: peta yang dapat ditelusuri kembali ke sumber peraturan dan data eksisting.
	\end{itemize}
\end{frame}

% Slide 16
\begin{frame}[mytitle=standard, t, mycolor=digiPH_gray]
	\frametitle{Studi Kasus dan Evaluasi Alternatif Rancangan}
	\begin{itemize}
		\item Gunakan satu tapak studi dengan data persil dan ketentuan lokal yang tersedia.
		\item Hitung kapasitas bangunan berdasarkan parameter regulasi yang telah diverifikasi.
		\item Bandingkan dua alternatif penempatan massa terhadap sempadan, akses, RTH, dan utilitas.
		\item Jelaskan kelebihan, keterbatasan, risiko, dan konsekuensi tiap alternatif.
		\item Rumuskan rekomendasi desain beserta alasan teknis dan dasar peraturannya.
	\end{itemize}
\end{frame}

% Slide 17
\begin{frame}[mytitle=standard, t, mycolor=white]
	\frametitle{Kesimpulan, Luaran, dan Referensi}
	\begin{itemize}
		\item Regulasi merupakan batas dasar pengambilan keputusan, sedangkan kondisi eksisting menentukan respons rancangan.
		\item Luaran mahasiswa: matriks ketentuan, perhitungan intensitas, peta kendala-potensi, dan rekomendasi penataan.
		\item Evaluasi berdasarkan ketepatan interpretasi, akurasi pemetaan, kelayakan rancangan, dan kualitas argumentasi.
		\item \textbf{Pertanyaan reflektif:} Bagaimana merancang tapak yang patuh regulasi tanpa mengabaikan konteks lingkungan dan kebutuhan pengguna?
	\end{itemize}
\end{frame}

\end{document}
